\cotefile{145}
\scanpage{20}
\mspage{8}

Il revient au $=$ de dire que $\sigma_0,\sigma_1,\sigma_\infty,\rho$ sont des générateurs, avec relations de définition

\[
(\overbrace{\rho^{-1}\sigma_\infty\sigma_0\deleted{\unsure{\rho^{-1}|\rho}}}^{l_\infty})
(\overbrace{\rho^{-1}\sigma_1\sigma_\infty\deleted{\unsure{\rho^{-1}|\rho}}}^{l_1})
(\overbrace{\rho^{-1}\sigma_0\sigma_1\deleted{\unsure{\rho^{-1}|\rho}}}^{l_0})=1 \tag*{(4)}
\]

\[
\sigma_0^2=\sigma_1^2=\sigma_\infty^2=1 \tag*{(5)}
\]

\[
\rho^3=1 \tag*{(6)}
\]

% slanted
\margin{\uline{NB} $\sigma_0(l_0)=l_{\deleted{\unsure{0}}\infty}^{-1}l_1^{-1}$\\
$\sigma_0(l_1)=l_\infty$\\
$\sigma_0(l_\infty)=l_1\deleted{\unsure{l_0^{-1}}}$\\
et itou pour $\sigma_1,\sigma_\infty$}

% slanted
\margin{Il faut ajouter la relation\unsure{,} $\rho\sigma_0\rho^{-1}=\sigma_1$ \deleted{etc} On a deux gén. $\sigma_0,\rho$ en fait avec relations $\sigma_0^2=\rho^3=1$, la \deleted{$[\unsure{\underbrace{\rho^{-1}\sigma_0\rho}_{\sigma_\infty}\sigma_0\rho^{-1}\underbrace{\rho\sigma_0\rho^{-1}}_{\sigma_1}\underbrace{\rho^{-1}\sigma_0\rho}_{\sigma_\infty}\rho^{-1}}]$} relation (4) étant conséquence des \unsure{précédentes|précédtes} (quand on y fait $\sigma_1=\rho\sigma_0\rho^{-1}$, $\sigma_\infty=\rho^{-1}\sigma_0\rho$). Noter que\\
$l_0=\deleted{\unsure{(\sigma_0\rho)^2}}\rho^{-1}(\sigma_0\rho)^2\rho=\boxed{\unsure{\rho^{-1}\sigma_0\rho\sigma_0\rho^{-1}}}$\\
$l_1=(\deleted{\unsure{\rho}}\sigma_0\rho)^2$\\
$l_\infty=(\rho^{-1}\sigma_0\deleted{\unsure{\rho}})^2$\\
$\mathrm{int}(\sigma_0)\,l_1=l_\infty$\\
$\mathrm{int}(\sigma_0)\,l_\infty=l_1$\\
$\mathrm{int}(\sigma_0)(l_0)=l_\infty^{-1}l_1^{-1}$}

\uline{NB} Si on y fait \inserted{de plus} $\sigma_0,\sigma_1,\sigma_\infty=1$, on trouve comme relation de définition $\rho^3=1$, \deleted{\unsure{donc}} sans plus, ce qui prouve que $\rho$ n'est pas un mot p.r. aux $\sigma_i$ \ldots

Tout ceci est compatible avec les notations du début, \deleted{car dans ces dernières on avait tout pris au \unsure{premiers|premier}, de \unsure{signes|gén.} \unsure{intervertis|interposés} $l'_i=l_i^{-1}$, $\sigma'_i=\sigma_{\dot\rho(i)}$, \unsure{donc} \uline{\unsure{remplacer}} $\sigma'$} \deleted{\unsure{puis} $\sigma_i$} si on y remplace les $l_i$ par les $l_i^{-1}$. Ici on a, comme $\sigma_i=\sigma'_{\unsure{\dot\rho^{-1}(i)|\rho^{-1}(i)|\dot\rho(i)}}$

\[
\left\{\begin{aligned}
\mathrm{int}(\sigma_i)&:\ \text{échange } l_j \text{ et }l_{\deleted{\unsure{j}}k}
 \ \inserted{(i,j,k)\ \text{distincts}}\quad\deleted{\text{et}}\\[-2pt]
&\deleted{\text{(et sa valeur sur }l_{\unsure{\dot\rho^{-1}(i)|\rho^{-1}(i)|\dot\rho(i)}}\text{ s'en déduit par (1))}}\\
\mathrm{int}(\rho)\,l_i&=l_{\dot\rho(i)}
\end{aligned}\right. \tag*{(7)}
\]

On voudrait maintenant expliciter les $\varepsilon_i$ avec $\varepsilon_i^2=l_i$ $(i=0,1,\infty)$, $\varepsilon_i$ étant donc un générateur \unsure{du|des} centralisateur (ou normalisateur, c'est itou) des $\unsure{L_i|l_i|b_i}$ dans ${\unsure{E|\Sigma|\mathscr E}}^+$. (ou \unsure{le centralisateur|les centralisateurs} dans $E$). Les calculs faits plus haut, avec les chgts de notation indiqués, donnent

\[
\left\{\begin{aligned}
\varepsilon_0&=l_0^{\unsure{1|-1}}\sigma_1l_{\unsure{1}}^{\unsure{1|-1}}\\
\varepsilon_1&=l_\infty^{-1}\sigma_\infty l_1^{-1}\\
\varepsilon_\infty&=l_0^{-1}\sigma_0 l_\infty^{-1}
\end{aligned}\right. \tag*{\deleted{\unsure{(8)}}}
\]
% archive sheet number: 19
\scanpage{21}

Donc on a les relations

\[
\left\{\begin{aligned}
\widetilde\sigma_0^2&=\widetilde\sigma_1^2=\widetilde\sigma_\infty^2=1\\
\widetilde\sigma_0\widetilde\sigma_1&=\widetilde\sigma_1\widetilde\sigma_\infty=\widetilde\sigma_\infty\widetilde\sigma_0=\rho^{\deleted{\unsure{\sim}}}\deleted{\unsure{(=\rho)}}\\
(\rho^{\deleted{\unsure{\sim}}\,3}&=1\ \text{ est conséquence des précédentes})
\end{aligned}\right. \tag*{(13)}
\]

% slanted
\margin{NB $\rho\widetilde\sigma_i\rho^{-1}=\widetilde\sigma_{\dot\rho i}$ i.e\\
$\rho\widetilde\sigma_0\rho^{-1}=\widetilde\sigma_1$\\
$\rho\widetilde\sigma_1\rho^{-1}=\widetilde\sigma_\infty$\\
$\rho\widetilde\sigma_\infty\rho^{-1}=\widetilde\sigma_0$}

et l'image inverse $\unsure{\widetilde\Sigma|\widetilde E}$ de $\varphi({\unsure{\mathfrak S_3|\mathcal S_3}}^{\deleted{\unsure{\sim}}})\subset\Gamma$ (ss-groupe d'indice 2) est produit \unsure{1/2|s/d} direct de $\mathfrak S_3$ et de $\pi$. Il faut expliciter les opérations des $\widetilde\sigma_i$ sur $\pi$ pour avoir la structure complète du groupe, On trouve que $\widetilde\sigma_{\unsure{i|0}}$ est au-dessus de $\varphi(\dot\sigma_{\unsure{i|0}}^{\deleted{\unsure{\sim}}})$ i.e. la symétrie p.r. au plan passant par $P,\overline P,\deleted{\unsure{\infty}}\,i$, qui transforme $l_0$ en $l_0^{-1}$, $\deleted{\unsure{l}}l_\infty$ en $l_{\unsure{1}}^{-1}$, $l_1$ en $l_\infty^{-1}$ :

% slanted
\margin{Ce sont \unsure{elles|celles} \unsure{se trouvent par|obtenues par} transport de structure, via les opérations de $\mathfrak S_3$ sur $(\unsure{U|V},P)$}

% slanted
\margin{\uline{Critique} La notation $\widetilde\sigma_0$ suggère \unsure{fautivt} que $\widetilde\sigma_0$ est au-dessus de $\sigma_0$, -- il vaut mieux écrire $\dot\sigma'_i=\dot\sigma_i\tau$ et noter $\widetilde\sigma'_i$ au lieu des $\widetilde\sigma_i$ $(i\in\{0,1,\infty\})$, ou mieux $\sigma'_i$ (sans $\sim$) -- pour réserver $\widetilde\sigma'_i$ à des éléments de $S\Sigma$, cf plus bas \ldots}

\[
\left\{\begin{aligned}
\widetilde\sigma_0 l_0\widetilde\sigma_0^{-1}&=l_0^{-1}&
\widetilde\sigma_0 l_1\widetilde\sigma_0^{-1}&=l_\infty^{-1}&
\widetilde\sigma_0 l_\infty\widetilde\sigma_0^{-1}&=l_1^{-1}\\
\widetilde\sigma_1 l_1\widetilde\sigma_1^{-1}&=l_1^{-1}&\ldots\\
&\ldots
\end{aligned}\right. \tag*{(14)}
\]

Les relations (13), (14), plus la relation fondamentale $l_\infty\inserted{\unsure{\sim|\checkmark}}l_1l_0=1$ (1), font une présentation de $\unsure{\widetilde G|\widetilde\Gamma}$, dont on peut si on veut éliminer $\rho$, en omettant d'écrire la relation qui le \unsure{définit|définist}\unsure{)|.}

On peut maintenant remonter $\dot\tau$ en trois manières, \unsure{à|6} \unsure{par|peu}

\[
\left\{\begin{aligned}
\tau_0&=\sigma_0\widetilde\sigma_0=\widetilde\sigma_0\sigma_0\\
\tau_1&=\sigma_1\widetilde\sigma_1=\widetilde\sigma_1\sigma_1\\
\tau_\infty&=\sigma_\infty\widetilde\sigma_\infty=\widetilde\sigma_\infty\sigma_\infty
\end{aligned}\right. \tag*{(15)}
\]

\deleted{donc}

\[
\left\{\begin{aligned}
\tau_0^2&=1\\
\tau_1^2&=1\\
\tau_\infty^2&=1
\end{aligned}\right. \tag*{(16)}
\]
% archive sheet number: 20
\scanpage{22}
\mspage{9}

\[
\left\{\begin{aligned}
\rho^3&=1\\
\varepsilon_\infty^2\,\varepsilon_1^2\,\varepsilon_0^2&=1\\
\deleted{\unsure{\varepsilon_0}}\ \varepsilon_1&=\varepsilon_0^{\unsure{\deleted{4}|4}}\varepsilon_\infty^4\rho\varepsilon_0^2\deleted{\unsure{\rho}}\varepsilon_\infty^2\\
\varepsilon_\infty&=\varepsilon_1\varepsilon_0^4\rho\varepsilon_1^2\deleted{\unsure{\rho}}\varepsilon_0^2\\
\varepsilon_0&=\varepsilon_\infty\varepsilon_1^4\rho\varepsilon_\infty^2\deleted{\unsure{\rho}}\varepsilon_1^2
\end{aligned}\right. \tag*{(10)}
\]

% slanted
\margin{il faut aussi ajouter (au moins!) les relations $\rho\varepsilon_0\rho^{-1}=\varepsilon_1$ etc donc on a deux gén. $(\varepsilon_0,\rho)$, \unsure{ça}}

On peut encore se demander si $\rho$ s'exprime en fonction des $\varepsilon_i$, et on trouve comme tantôt que non.

Quant à relever $\dot\tau$ en un élément de $\unsure{\Sigma|E}$, \deleted{par} un chemin de $P$ à $\overline P$, on ne voit pas de choix privilégié -- on voit trois choix également naturels. Au lieu de les expliciter \unsure{directmt|d'emblée}, \unsure{je préfère|je puis faire} considérer le sous-groupe d'isotropie de $P$ dans $\Gamma=\Gamma_{0,3}=\unsure{\mathfrak S_3|\mathcal S_3}\times\mathbb Z/2$, qui est formé des couples $\dot\sigma,\dot\tau^\varepsilon$ tels que $\mathrm{sg}(\dot\sigma)=\varepsilon$ ; \deleted{On} i.e. les éléments

\[
\widetilde{\dot\sigma}=(\dot\sigma,\dot\tau^{\mathrm{sg}(\dot\sigma)})\in\Gamma
\]

qui sont relevés canoniquement en des éléments $(\widetilde{\dot\sigma},\text{\unsure{lacet} trivial en }P)$ de $\Sigma$. On trouve une \unsure{Imm|imm|hom}

\[
\left\{\begin{aligned}
\varphi:\mathfrak S_3&\longrightarrow\Gamma=\mathfrak S_3\times\mathbb Z/2\\
\dot\sigma&\longmapsto\dot\sigma,\mathrm{sg}(\dot\sigma)
\end{aligned}\right. \tag*{(11)}
\]

et un relèvement canonique de $\varphi$ en

\[
\sigma\longrightarrow\widetilde\sigma:\quad\mathfrak S_3\longrightarrow\Sigma \tag*{(12)}
\]

On pose $\widetilde\sigma_i=\widetilde{\dot\sigma_i}$ \inserted{(on a d'autre part} $\deleted{\unsure{\rho_0}}\ \widetilde\rho=\rho$
% archive sheet number: 21
\scanpage{23}
\mspage{10}

Le sous-groupe $\unsure{\Sigma|E}_\tau$ de $\Sigma$ image inverse du $\mathbb Z/2$ engendré par $\tau$, est le groupe cartographique ``\unsure{tri-angulé|triangulé} pondéré'' de Malgoire, ayant comme \deleted{\unsure{relation}} \inserted{gén. $\tau_0,\tau_1,\tau_\infty$} avec

\[
\left\{\begin{aligned}
\deleted{\tau_1\tau_\infty={}}l_0&=\tau_{\deleted{\unsure{0}}\infty}\tau_{\deleted{\unsure{\infty}}1}\\
\deleted{\tau_\infty\tau_0={}}l_1&=\tau_{\deleted{\unsure{1}}0}\tau_\infty\\
l_\infty&=\tau_1\tau_0
\end{aligned}\right. \tag*{(17)}
\]

\margin{\uline{NB} --}

$\Sigma_\tau$ peut être considéré comme le groupe ayant comme \inserted{gén. les $\tau_i$, avec les} seules relations (16) -- on a un hom signature évident $\tau_i\longrightarrow-1$ de $\Sigma_\tau$ dans $\pm1$, dont le noyau est $\pi$ engendré par les $l_i$ décrits par (17) \ldots.

\centerline{\rule{1.2cm}{0.4pt}}

% slanted
\margin{On a les relations \unsure{$\rho$}}

Revenons à une description plus \unsure{symétrique|symbolique} des autom. extérieurs à lacets de $\Sigma$. On désigne par $\widehat\Sigma$ le complété profini de $\Sigma$, extension de $\widehat{\unsure{\mathscr T|\Gamma}}_{0,3}=\mathscr T_{03}=\unsure{\mathfrak S_3|\mathcal S_3}\times\mathbb Z/2$ par $\widehat\pi$, et par $\widehat{\widehat\Sigma}$ le groupe des autom. \deleted{extérieurs} lacets de $\widehat\pi$, extension de $\widehat{\widehat\Gamma}=\mathrm{Aut\,ext}_{\mathrm{lac}}(\widehat\pi)$ par $\widehat\pi$

\diagram{\begin{tikzcd}[ampersand replacement=\&,column sep=small,row sep=0.2em]
\text{(18)} \& 1\arrow[r] \& \widehat\pi\arrow[r] \& \widehat{\widehat\Sigma}\arrow[r] \& \widehat{\widehat\Gamma}\arrow[r]\arrow[d,equal] \& 1\\
\& \& \& \& \mathrm{Aut\,ext}_{\mathrm{lac}}(\widehat\pi) \&
\end{tikzcd}}
% archive sheet number: 22
\scanpage{24}

qui contient $\widehat{\unsure{\Sigma|E}}$ comme sous-extension. On désigne par $\mathbb{\Gamma}$ le commutant de $\unsure{\mathfrak S_3|\mathcal S_3}$ dans $\widehat{\widehat\Gamma}$, \inserted{(qui contient $\dot\tau$)} et par $\mathbb E$ l'extension correspondante de $\mathbb{\Gamma}$ par $\widehat\pi$. Notons déjà que

\[
\left\{\begin{aligned}
\mathbb{\Gamma}\cap\Gamma\deleted{\unsure{\mathfrak S_3}}&=\{1,\tau\}\\
\mathbb{\Gamma}\cap\mathfrak S_3&=\{1\}
\end{aligned}\right.\qquad\deleted{\unsure{\mathrm{NB}}}
\]

% slanted
\margin{\uline{NB} \diagram{$\begin{array}{ccc}&\mathfrak S_3\times\widehat{\widehat\Gamma}^{!}&\\&\rotatebox{90}{$\subset$}&\\&\widehat{\widehat\Gamma}\supset\deleted{\unsure{\Gamma}}\mathfrak S_3\times\mathbb{\Gamma}&\end{array}$}
-- j'ignore si cette inclusion est une égalité, i.e. si $\mathbb{\Gamma}\subset\widehat{\widehat\Gamma}^{!}$ est une égalité}

On veut donner d'abord une description plus symétrique des éléments de $\widehat{\widehat\Sigma}^{!}$, sans faire jouer un rôle particulier à $L_0$ engendré par $l_0$ $(L_0=l_0^{\widehat{\mathbb Z}})$. Un $u\in\unsure{\widehat\Sigma|\widehat{\widehat\Sigma}}$ est \inserted{peut être} défini par un multiplicateur

\[
\unsure{p|\mathfrak p}=\chi(u)\in\widehat{\mathbb Z}^{*} \tag*{(19)}
\]

et des

\[
g_i\in\widehat{\unsure{\deleted{G}\pi|\pi}}\qquad i=0,1,\infty \tag*{(20)}
\]

\unsure{par|pour}

\[
u(l_i)=\mathrm{int}(g_i^{\unsure{\deleted{-1}|-1}})\,l_i^p \tag*{(21)}
\]

les données (19), (20) étant liées par la relation

\[
\bigl(\mathrm{int}(g_\infty^{\unsure{\deleted{-1}|-1}})\,l_\infty^p\bigr)
\bigl(\mathrm{int}\,g_1^{\deleted{-1}}\cdot l_1^p\bigr)
\bigl(\mathrm{int}(g_0^{\deleted{-1}})\,l_0^p\bigr)=1 \tag*{(22)}
\]

et par le fait que les $u(l_i)$ définis par (21) engendrent $\widehat\pi$ topologiquement. Pour $u$ fixé\unsure{,} $p$
% archive sheet number: 23
\scanpage{25}
\mspage{11}

est unique, et $g_i$ défini multiplication à \deleted{\unsure{droite}}\unsure{gauche} par un él. quelconque de $L_i$, soit $l_i^{\alpha_i}$ $(\alpha_i\in\widehat{\mathbb Z})$. On passe de $u$ à $\mathrm{int}(g)\circ u$ en \inserted{gardant $p$, et en} remplaçant les $g_i$ par les $gg_i$.

Si $v$ est donné par $q,(h_i)$, on a

\[
vu(l_i)=v(\mathrm{int}(g_i^{\unsure{\deleted{-1}|-1}})\,l_i^p)=\mathrm{int}(v(g_i^{\unsure{\deleted{-1}|-1}}))\,v(l_i)^p
\]

\[
\text{où}\quad v(l_i)=\mathrm{int}(h_i^{\deleted{-1}})\,l_i^q\qquad\text{donc}\quad
v(l_i)^p=\mathrm{int}(h_i^{\deleted{-1}})\,l_i^{pq}
\]

donc

\[
vu(l_i)=\mathrm{int}\bigl(\overbrace{v(g_i^{\deleted{-1}})\,h_i^{\deleted{-1}}}^{\deleted{\unsure{(h_i\,v(g_i))^{-1}}}}\bigr)\,l_i^{pq} \tag*{(23)}
\]

Si on désigne par $S\widehat{\widehat{\unsure{\Sigma|E}}}^{!}$ l'extension de $\widehat{\widehat\Sigma}^{!}$ par $\widehat{\mathbb Z}^{\{0,1,\infty\}}$ définie \deleted{par, \unsure{le groupe}} groupe des \deleted{couples} tuples d'un $u\in\widehat{\widehat\Sigma}^{!}$ et d'un système de $g_i$ qui explicitent \inserted{via (21)} l'\unsure{invariance|isomorphisme} de la structure à lacets ($p$ \unsure{est} le multiplicateur), on a sur $S\widehat{\widehat\Sigma}^{!}$ trois fonctions à valeurs ds $\deleted{\unsure{G}}\widehat\pi$

% slanted
\margin{\uline{NB} $S\widehat{\widehat\Sigma}^{!}$ a été construit comme le groupe des systèmes $(p,g_0,g_1,g_\infty)\in\widehat{\mathbb Z}^{*}\times\pi^{[0,1,\infty]}$ satisfaisant la relation (22)}

\[
(u,(g_i))\overset{\gamma_i}{\longmapsto}g_i,\qquad\text{et on aura}
\]

\[
\gamma_i(\widetilde v\widetilde u)=
\deleted{\widetilde v(\unsure{u_i|\gamma_i}(\widetilde u))\cdot\gamma_i(\widetilde v)}
\deleted{\unsure{\gamma_i(\widetilde v)\cdot v(\gamma_i(\widetilde u))}}
\inserted{\widetilde v(\gamma_i(\widetilde u))\gamma_i(\widetilde v)} \tag*{(24)}
\]

i.e. $(\gamma_i)$ est un 1-cocycle sur le groupe \deleted{\unsure{opposé}} \inserted{$S\widehat{\widehat\Sigma}^{!}$} à valeurs dans le groupe \deleted{\unsure{opposé au groupe}} \inserted{opposé au groupe} $\pi^{\{0,1,\infty\}}$.

\deleted{\unsure{groupe}} Le sous-groupe de $S\widehat{\widehat\Sigma}^{!}$ noyau de $S\widehat{\widehat\Sigma}\longrightarrow\mathrm{Aut\,ext}_{\mathrm{lac}}^{!}(\widehat\pi)$ \deleted{\unsure{s'i}} s'identifie au groupe produit $\bigl(\widehat\pi\times\prod_{i\in\{0,1,\infty\}}L_i\bigr)\simeq\widehat\pi\times\deleted{\unsure{\widehat\pi}}\widehat{\mathbb Z}^{\{0,1,\infty\}}$
% archive sheet number: 24
\scanpage{26}

en associant à $(g,(\alpha_i)_{i\in\{0,1,\infty\}})$ le système \deleted{des} $p=1$, $(g_i)=\deleted{\unsure{(l_i^{\alpha_i})}}\deleted{\unsure{(l_i^{\alpha_i}g^{-1})}}(g\,l_i^{\alpha_i})$

% slanted
\margin{\uline{NB} Il y a lieu d'introduire un groupe $S\widehat{\unsure{\Sigma|E}}$ extension de $\Sigma$ par $\mathbb Z^{\{0,1,\infty\}}$ (déduite d'une extension canonique de $\Gamma$ par $\mathbb Z^{\{0,1,\infty\}}$) et \unsure{d'épingler|d'expliciter|d'introduire} les relèvements des \deleted{\unsure{de}} $\sigma_i,\rho$ etc en des éléments \inserted{déjà} $\widetilde\sigma_i,\widetilde\rho$ de $S\Sigma$ \ldots.}

Il faudrait expliciter les opérations des éléments $\sigma_i$ ou $\varepsilon_i$ \deleted{sur} $(\in S\widehat{\widehat\Sigma}^{!})$ sur $\deleted{\unsure{\Sigma}}\widetilde u\in S\widehat{\widehat\Sigma}^{!}$ \deleted{\unsure{de base}} par conjugaison, et en particulier l'action de $\rho$. On trouve que \deleted{$\chi$ \unsure{(bien sûr)}} les $\gamma_i$ sont ``covariants'' par $\rho$, i.e.

\[
\gamma_{\unsure{i|i'}}(\rho\widetilde u\rho^{-1})=\gamma_{i+1}\unsure{(\bar u)|(\widetilde u)}
\]

\deleted{i.e. si $\widetilde u$ est défini par}
% archive sheet number: 25
\scanpage{27}
\mspage{12}

L'extension $\unsure{\Sigma|E}=\Sigma_{0,3}$ de $\Gamma=\Gamma_{03}=\underbrace{\unsure{\mathfrak S|\mathcal S}_{[0,1,\infty]}}_{\text{ou }\mathfrak S_3}\times\mathbb Z/2\mathbb Z$ \inserted{engendré par $\tau:z\mapsto\overline z$} par $\pi=\pi_{0,3}=\{l_0,l_1,l_\infty\mid l_\infty l_1l_0=1\}=\pi_1(\underbrace{\mathbb P^1_{\unsure{\mathbb C|C}}\setminus\{0,1,\infty\}}_{U=U_{0,3}},-\bar j\}$.

Les classes de conjugaison de scindages partiels maximaux \deleted{au-dessus de} correspondent aux points fixes de $\dot\rho,\dot\sigma_0,\dot\sigma_1,\dot\sigma_\infty$ opérant sur $X_{0,3}$, et \inserted{(éventuellt)} aux comp. connexes des ens. d'invariants de $\deleted{\unsure{\rho}}\tau,\dot\sigma_0\tau,\dot\sigma_1\tau,\dot\sigma_\infty\tau$. Les pts fixes de $\rho$ sont $P$ et $\overline P$ dont le stabilisateur est d'ordre 6, \inserted{savoir $\varphi(\mathfrak S_3)$} formé de $1,\dot\rho,\dot\rho^2,\dot\sigma_1\tau,\dot\sigma_2\tau,\dot\sigma_3\tau$.

On a étudié le relèvement correspondant \inserted{à $P$} de $\varphi(\mathfrak S_3)$, on trouve des relèvements ($\pi$-conjugués \inserted{entre eux}) relatifs à $\overline P$ s'obtiennent en \deleted{\unsure{conj}} prenant un élément de $\Sigma$ \deleted{\unsure{q}} qui \deleted{\unsure{soit|sont}} au dessus d'un élément de $\Gamma$ envoyant $P$ en $\overline P$ (savoir $\dot\sigma_0,\dot\sigma_1,\dot\sigma_\infty,\tau,\tau\rho,\tau\rho^2$) et en conjuguant par cet élément

on peut donc prendre $\deleted{\unsure{\sigma_0,\tau}}$ un des éléments $\sigma_0,\sigma_1,\sigma_\infty$ (ou $\varepsilon_0,\varepsilon_1,\varepsilon_\infty$) ou $\tau_0,\tau_1,\tau_\infty$ ou \deleted{\unsure{les}} les produits des précédents par \deleted{\unsure{un}} $\rho$ ou $\rho^2$\ldots.

\deleted{Les} Le \inserted{seul} pt\deleted{s} fixe\deleted{s} de $\sigma_0$ est 2, dont le stabilisateur est le groupe $\simeq\mathbb Z/2\times\mathbb Z/2$ engendré par $\sigma_0$ et $\tau$. On a un relèvt de ce sous-groupe
% archive sheet number: 26
\scanpage{28}

dans $\unsure{\Sigma|E}$, qu'on peut choisir canoniquement en prenant le quart de grand cercle qui joint $P$ au milieu $2$ du \deleted{segment} l'arc $(\unsure{1|0},\infty)$ de grand cercle \unsure{opposée|opposé} à $0$, ce qui définit un isom du rev universel de $\deleted{\unsure{U}}U$ en $P$, avec celui de $U$ en $\bar P$ sur le quel le groupe précédent opère de façon canonique, donc aussi par transport de structure sur $\widetilde U\deleted{\unsure{(\bar P)}}(P)$, donc il s'envoie dans $\Sigma$.

On trouve que

\[
\dot\sigma_0\mapsto\sigma_0,\quad\dot\tau\longrightarrow\tau_0
\]

est le remontage correspondant. De $=$ pour le pt fixe $-1$ de $\dot\sigma_{\unsure{1}}$, ou $\frac12$ de $\dot\sigma_\infty$, on trouve des scindages partiels sur les ss-groupes de \deleted{\unsure{conj}} stabilité $\Gamma_{-1}\subset\Gamma$ et $\Gamma_2\subset\Gamma$. Il est commode de poser

\[
\deleted{\unsure{P}}\quad 2=0',\quad-1=1',\quad1/2=\infty',
\]

d'où des groupes d'isotropie $\Gamma_{0'}=\{\dot\sigma_0,\tau\}$, $\Gamma_{1'}=\{\dot\sigma_1,\tau\}$, et $\Gamma_{\infty'}=\{\dot\sigma_\infty,\tau\}$, qui se remontent resp en $\sigma_0$ et $\tau_0$, \deleted{\unsure{ou}} $\sigma_1$ et $\tau_1$, en $\sigma_\infty$ et $\tau_\infty$.

Les composantes connexes de $U^\tau$ sont les trois arcs de grand cercle entre $0,1,\infty$, \deleted{qui} \unsure{contient|contenant} chacun un \deleted{pt} et un seul des pts $0',1',\infty'$, et les classes de conjugaison
% archive sheet number: 27
\scanpage{29}
\mspage{13}

correspondantes de relèvements \inserted{de $\{1,\dot\tau\}$} sont resp. les relèv\unsure{ts} en $\tau_0$, en $\tau$, et en $\tau_\infty$ -- qui ne donnent rien de nouveau. Reste à considérer $\tau\dot\sigma_1$ (et ses conjugués \inserted{par $\rho$} $\tau\dot\sigma_\infty,\tau\unsure{\dot\sigma_0}$), on a $\dot\sigma_1(2)=1/2$, donc

\[
(\tau\dot\sigma_1)(2)=\frac12
\]

\[
\text{et}\qquad U^{\tau\dot\sigma_1}=\{z\in U\mid z\bar z=1\}=\mathbb U-\{1\},
\]

où $\mathbb U=\{z\in\mathbb C\mid |z|=1\}$. Cet ensemble contient $P$ et $\bar P$, et ne donne donc pas de nouveau relèvement p.r. à celui déjà donné, avec $\widetilde\sigma_1$ relevant $\dot\sigma_1\tau$.

Il faudrait quand-$=$ donner aussi une description de $\unsure{\Sigma|E}$ tt entier par gén. et relations, en partant d'un syst. de gén. de $\Gamma=\unsure{\mathfrak S_3|\mathcal S_3}\times\mathbb Z/2$, qu'on va relever de façon judicieuse. On prendra les \inserted{quatre} générateurs

\deleted{$\rho'=\rho\tau,\unsure{\tau}\widetilde{\dot\sigma}_0,\widetilde{\dot\sigma}_1,\widetilde{\dot\sigma}_\infty$, avec les relations}

\diagram{\deleted{$\begin{array}{ccc}\Vert&\Vert&\Vert\\\dot\sigma_0\tau&\dot\sigma_1\tau&\dot\sigma_\infty\tau\end{array}$}}

\deleted{(donc $\rho'{}^6=1,\rho'{}^3=\tau,\rho'{}^2=\rho^2=\rho^{-1}$)}

\[
\left\{\begin{aligned}
&\deleted{\widetilde{\dot\sigma}_1^2=\widetilde{\dot\sigma}_{\unsure{2|1}}^2=\widetilde{\dot\sigma}_3^2=1}\\
&\deleted{\widetilde{\dot\sigma}_0\widetilde{\dot\sigma}_1=\widetilde{\dot\sigma}_1\widetilde{\dot\sigma}_\infty=\widetilde{\dot\sigma}_\infty\widetilde{\dot\sigma}_0=\dot\rho^{\unsure{-2|2}}}\\
&\deleted{\rho'{}^6=1}\\
&\deleted{\rho'{}^3\ \text{commute aux }\dot\sigma_i}
\end{aligned}\right.
\]
% archive sheet number: 28
\scanpage{30}

\deleted{qui décrivent $\unsure{\mathfrak S_3|\mathcal S_3}\times\mathbb Z/2$, on va relever ces générateurs en des éléments de $\unsure{\Sigma|E}$, de façon que $l_\infty,l_1,l_0$ s'expriment comme des mots en ces éléments, et il suffira alors d'ajouter cette relation $l_\infty l_1l_0$. aux relations correspondant aux précédentes \ldots}

$\dot\sigma_0,\dot\sigma_1,\dot\sigma_\infty,\deleted{\unsure{\rho}}\dot\rho,\tau$, avec les relations

\[
\begin{aligned}\bullet\\\bullet\\\bullet\\\bullet\\\bullet\end{aligned}\,
\left\{\begin{aligned}
\dot\sigma_0^2&=\dot\sigma_1^2=\dot\sigma_\infty^2=1\\
\dot\sigma_0\dot\sigma_1&=\dot\sigma_1\dot\sigma_\infty=\dot\sigma_\infty\dot\sigma_0=\dot\rho\\
\dot\rho^3&=1\\
\tau^2&=1\\
\tau&\ \text{commute à }\dot\sigma_0,\dot\sigma_1,\dot\sigma_3\quad\text{(donc aussi à }\dot\rho\text{)}
\end{aligned}\right.
\]

on les relèvera en $\deleted{\tau_0,\tau_1,\tau_\infty}\ \sigma_0,\sigma_1,\sigma_\infty,\rho,\underbrace{\tau_0,\tau_1,\tau_\infty}_{\text{sur }\tau}$ avec les relations

% slanted
\margin{\deleted{NB} Il faut ajouter les relations $\left\{\begin{aligned}\rho\sigma_0\rho^{-1}&=\sigma_1\\\rho\sigma_1\rho^{-1}&=\sigma_\infty\\\rho\sigma_\infty\rho^{-1}&=\sigma_0\end{aligned}\right.$ \deleted{\unsure{sont}} conséquence\unsure{s} des relations ci-contre}

\margin{[ NB \unsure{On posera|On trouve} $l_0=\tau_\infty\tau_1$, $\deleted{l_1=\tau_0\tau_\infty}$, $l_1=\tau_0\tau_\infty$ $\deleted{\unsure{l_\infty={}}}$ $l_\infty=\tau_1\tau_0$ et $l_\infty l_1l_0=1$ automatique ]\\
$\uparrow$ \deleted{\unsure{les}} Ceci implique que $\tau_0,\tau_1,\tau_\infty$ engendrent \unsure{un ss-groupe engendré|le groupe contenant} par $l_0,l_1,l_\infty$}

\[
\left\{\begin{aligned}
&\sigma_0^2=\sigma_1^2=\sigma_\infty^2=\tau_0^2=[\tau_1^2=\tau_\infty^2]=1\\
&\rho^3=1\\
&\rho\tau_0\rho^{-1}=\tau_1,\quad\rho\tau_1\rho^{-1}=\tau_\infty,\quad\rho\tau_\infty\rho^{-1}=\tau_0\\
&\tau_0\sigma_{\deleted{\unsure{\infty}}0}=\sigma_0\tau_0,\quad\tau_1\sigma_1=\sigma_1\tau_1,\quad\tau_\infty\sigma_\infty=\sigma_\infty\tau_\infty\\
&\sigma_0\sigma_1=\tau_\infty\tau_1\rho,\quad\sigma_1\sigma_\infty=\tau_0\tau_\infty\rho,\quad\sigma_\infty\sigma_0=\tau_1\tau_0\rho
\end{aligned}\right.
\]

\margin{Ce sont les éléments $\widetilde\sigma_0,\widetilde\sigma_1,\widetilde\sigma_\infty$ $\longrightarrow$}

\deleted{\uline{Question} \deleted{\unsure{Peut-}} \inserted{On peut} exprimer $\rho$ en fonction des $\sigma_0,\sigma_1,\sigma_\infty,\tau_0,\unsure{1|\tau_1},\tau_\infty$\unsure{??}, \deleted{et non plus} \inserted{par l'une du} dernier groupe de relations, et substituer dans les autres.}
% archive sheet number: 29
\scanpage{31}
\mspage{14}

On peut prendre aussi comme seuls gén. $\rho,\sigma_0,\tau_0$ (dont se déduisent $\sigma_1,\tau_1$, \inserted{$\sigma_{\unsure{0|\infty}},\tau_\infty$,} comme $\rho\sigma_0\rho^{-1},\rho\tau_0\rho^{-1},\rho^2\sigma_0\rho^{-2},\rho^2\tau_0\rho^{-2}$) et les relations

\[
\left\{\begin{aligned}
\text{\unsure{a}) }&\sigma_0^2=\tau_0^2=1\quad \tau_0\sigma_0=\sigma_0\tau_0\quad
\text{(i.e. }(\overbrace{\sigma_0\tau_0}^{\widetilde\sigma_0})^2=1\text{)}\\
\text{b) }&\rho^3=1\\
&\sigma_0\underbrace{\rho\sigma_0^{-1}\rho^{-1}}_{\sigma_1}
=\underbrace{\rho^2\tau_0\rho^{-2}}_{\tau_\infty}\underbrace{\rho\tau_0\rho^{-1}}_{\tau_1}\rho\quad
\text{i.e. }\ \sigma_0\rho\sigma_0^{-1}\rho^{-1}=\rho^{-1}\tau_0\rho\tau_0\\
&\text{ou enfin}\\
&\rho\sigma_0\rho^{-1}\sigma_0=\tau_0\rho^{-1}\tau_0\rho\quad
\text{soit }\ \deleted{\text{\unsure{i.e. ou enco.}}}\ \deleted{\mathrm{int}(\rho)\sigma_0}\ 
[\rho,\sigma_0]=[\tau_0,\rho^{-1}]\\
&\text{soit }\quad [\rho,\sigma_0][\rho^{-1},\tau_0]=1
\end{aligned}\right.
\]

C'est donc un quotient du produit libre $(\mathbb Z/2\mathbb Z\times\mathbb Z/2\mathbb Z)*\mathbb Z/3\mathbb Z\unsure{)}$ par le sous-groupe invariant décrit par la relation\unsure{s} $[\rho,\sigma_0][\rho^{-1},\tau_0]=1$.

\uline{NB} On trouve une symétrie dans le groupe, qui échange $\sigma_0$ et $\tau_0$, et échange $\rho$ et $\rho^{-1}$. [ donc qui échange aussi $\sigma_1$ et $\tau_\infty$, $\sigma_\infty$ et $\tau_1$ ] -- mais cette symétrie n'est pas compatible avec la structure d'extension, comme on voit tout \deleted{\unsure{de s}} de suite (car en passant au quotient $\unsure{\mathfrak S_3|\mathcal S_3}\times\mathbb Z/2$, elle devrait conserver \unsure{le centre|les centres} !).
% archive sheet number: 30
\scanpage{32}

Il faudrait encore expliciter l'extension $S\Gamma\unsure{.}=S\Gamma_{0,3}$ de $\Gamma$ par $\mathbb Z^{\{0,1,\infty\}}$, ce qui nous donnera (par produits fibrés) une description de $S\unsure{\Sigma|E}=S\Gamma\times_\Gamma\Sigma$.
% archive sheet number: 31
\scanpage{33}
\mspage{15}

\textcircled{2} \uline{Le groupe} $S\widehat{\widehat{\unsure{\Sigma|E}}}$.

Par le calcul, il s'explicite comme l'ens des systèmes

\[
(\underset{\substack{\rotatebox{-90}{$\in$}\\\widehat{\mathbb Z}^{*}}}{p},
 \underset{\substack{\rotatebox{-90}{$\in$}\\\unsure{\mathfrak S_3|\mathcal S_3}}}{\sigma},
 \underset{\substack{\rotatebox{-90}{$\in$}\\\widehat\pi}}{g_0},
 \underset{\substack{\rotatebox{-90}{$\in$}\\\widehat\pi}}{g_1},
 \underset{\substack{\rotatebox{-90}{$\in$}\\\widehat\pi}}{g_\infty})=\widetilde u \tag*{(1)}
\]

satisfaisant la condition

\[
\mathrm{int}\unsure{(g)_{\infty}|(g_{\infty})}(l_{\sigma(\infty)}^p)\cdot
\mathrm{int}(g_1)(l_{\sigma(1)}^p)\cdot
\mathrm{int}(g_\infty)(l_{\sigma(\unsure{\infty|0})}^p)=1 \tag*{(2)}
\]

et tels que les trois facteurs de ce produit soient un syst. de générateurs de $\widehat\pi$ [avec la relation génératrice (2)]. Le morphisme $u$ défini par ce \deleted{\unsure{morphi}} système est

\[
\left\{\begin{aligned}
u(l_0)&\unsure{\mapsto|=}\mathrm{int}(g_0)\,l_{\sigma(0)}^p\\
u(l_1)&=\mathrm{int}(g_1)\,l_{\sigma(1)}^p\\
u(l_\infty)&=\mathrm{int}(g_\infty)\,l_{\sigma(\infty)}^p
\end{aligned}\right. \tag*{(3)}
\]

ce qui explicite l'homom

\[
S\widehat{\widehat\Sigma}\xrightarrow{p}\widehat{\widehat\Sigma}=\mathrm{Aut}_{\mathrm{lac}}(\widehat\pi) \tag*{(4)}
\]

L'homomorphisme canonique

\[
i:\widehat\pi\times\widehat{\mathbb Z}^{\{0,1,\infty\}}\longrightarrow S\widehat{\widehat\Sigma} \tag*{(\unsure{5})}
\]

s'explicite par

% slanted
\margin{On \unsure{retrouve|retiendra} que (\unsure{4|p}) est surjectif, et a comme noyau $i(\widehat\pi)$}

\[
i(g;\alpha_0,\alpha_1,\alpha_\infty)=(1,1,g\,l_0^{\alpha_0},g\,l_{\unsure{0|1}}^{\alpha_1},g\,l_\infty^{\alpha_\infty})\quad\text{donc} \tag*{(6)}
\]

\[
\deleted{\text{donc}}\quad pi(g;\alpha_0,\alpha_1,\alpha_\infty)=\mathrm{int}\,g. \tag*{(7)}
\]

La multiplication dans $\widehat{\widehat\Sigma}$ est \unsure{donné|donnée} par

\[
\underbrace{(p,\sigma,g_0,g_1,g_\infty)}_{\widetilde u}
\underbrace{(p',\sigma',g'_0,g'_1,g'_\infty)}_{\widetilde u'}
=(pp',\sigma\sigma',u(g'_0)g_{\sigma'(0)},u(g'_1)g_{\sigma'(1)},u(g'_\infty)g_{\sigma'(\infty)}\unsure{)} \tag*{(8)}
\]

Cette formule n'est pas, du point de vue du calcul,
% archive sheet number: 32
\scanpage{34}

entièrement explicite et satisfaisante, puisque au deuxième membre interviennent les quantités $u(g'_0),u(g'_1),u(g'_\infty)$, qui ne sont donc pas exprimées directement en termes des \inserted{seules} données $(p,\sigma,g_0,g_1,g_\infty)$, $(p',\sigma',g'_0,g'_1,g'_\infty)$ ; \deleted{sauf} \inserted{mais ils s'explicitent ainsi} quand $g'_0,g'_1,g'_\infty$ sont eux-mêmes donnés explicitement comme des mots (ou limites de tels) en les $l_0,l_1,l_\infty$.

Cependant quand $\widetilde u$ est de la forme $i(g,\alpha_0,\alpha_1,\alpha_\infty)$, on trouve (en laissant tomber les $'$ du \unsure{$2^{\mathrm e}$|$2^{\mathrm d}$} facteur)

\[
 i(g,\alpha_0,\alpha_1,\alpha_\infty)(p,\sigma,g_0,g_1,g_\infty)=(p,\sigma,\deleted{\unsure{i(g|u(g}}gg_0l_{\deleted{\unsure{g_0|g}}\sigma0}^{\alpha_{\sigma(0)}},gg_1l_{\sigma1}^{\alpha_{\sigma(1)}},gg_\infty l_{\sigma\infty}^{\alpha_{\sigma\infty}})\tag*{(9)}
\]

qui se simplifie pour $\sigma=1$ en

\[
 i(g,\alpha_0,\alpha_1,\alpha_\infty)(p,1,g_0,g_1,g_\infty)=(p,\sigma,gg_0l_0^{\alpha_0},gg_1l_1^{\alpha_1},gg_\infty l_\infty^{\alpha_\infty})\tag*{($9'$)}
\]

Par ailleurs, l'opération de $\widehat{\widehat{\unsure{\Sigma|E}}}$ sur $i(\pi\times\mathbb Z^{\{0,1,\infty\}})$ par autom. int. s'explicite aisément, et \inserted{ne dépend que de $u$ \deleted{\unsure{et}},}

\[
 \widetilde u\,i(g;\alpha_0,\alpha_1,\alpha_\infty)\deleted{\unsure{=}}\widetilde u^{-1}=i(u(g);p\alpha_{\sigma^{-1}0},p\alpha_{\sigma^{-1}1},p\alpha_{\sigma^{-1}\infty})\tag*{(10)}
\]

où $\widetilde u=(p,\sigma,g_0,g_1,g_\infty)$

-- mais sur cette formule, on peut faire les mêmes commentaires que plus haut, sur le manque d'explicitation du \deleted{\unsure{terme}} $u(g)$ ds cette formule, pour $g$ non exprimé comme un mot p.r. aux $l_0,l_1,l_\infty$ (ou limite de tels mots)

% archive sheet number: 33
